\documentclass[11pt,a4paper]{article}
\usepackage[T1]{fontenc}
\usepackage[utf8]{inputenc}
\usepackage{amsmath,amssymb,amsthm}
\usepackage[margin=1in]{geometry}
\usepackage[colorlinks=true,linkcolor=blue,urlcolor=blue]{hyperref}
\theoremstyle{plain}
\newtheorem{theorem}{Theorem}
\newtheorem{lemma}[theorem]{Lemma}
\newtheorem{proposition}[theorem]{Proposition}
\newtheorem{corollary}[theorem]{Corollary}
\theoremstyle{definition}
\newtheorem{remark}[theorem]{Remark}

\title{Recovering the unwritten recurrences of the table OEIS A232076}
\author{Adrian Perez Fontelles\\ \small Independent researcher}
\date{2 September 2026}
\begin{document}
\maketitle

\begin{abstract}
OEIS A232076 is a two-dimensional table $T(n,k)$ whose empirical recurrences are, for several of
its lines, given only as an ORDER --- ``k=3: [order 10]'' and the like --- with the
coefficients in a linked file rather than in the entry. Those conjectures can still be settled,
and the recurrences recovered. A line of the table is a fixed-width or fixed-height array
count, hence a walk count on finitely many vertices, hence satisfies some monic recurrence of
order at most that number; Berlekamp--Massey on twice as many exact terms returns the MINIMAL
one. Where its order equals the order the entry states --- and where the same holds after the
first few terms are dropped --- any recurrence of that order the line satisfies has a
characteristic polynomial that is a multiple of the minimal one and of equal degree, hence is
that polynomial. This note settles column 3, column 4, row 2, row 3, row 4.
\end{abstract}

\noindent\small 2020 Mathematics Subject Classification. 05A15, 11B37, 15A18.\normalsize

\section{The table and the unwritten conjectures}

OEIS A232076 is ``T(n,k)=Number of (n+1)X(k+1) 0..1 arrays with every element equal to some horizontal, diagonal or antidiagonal neighbor, with top left element zero.''. It is read by antidiagonals and begins
\[
3, 15, 11, 46, 87, 34, 161, 520,\ \dots
\]
The entry carries, among its empirical claims:
\begin{quote}\small
k=3: [order 10]\\
k=4: [order 30]\\
n=2: [order 8]\\
n=3: [order 20]\\
n=4: [order 54]
\end{quote}
As of the ``Last modified'' line on the live entry (Jul 23 2025 06:52:15, revision 6) these are still
recorded as empirical, and the recurrences themselves are not in the entry text.

\section{Why a line of the table is a walk count}

Fix the index that the line fixes. The entry defines $T(n,k)$ as a number of arrays of a shape
determined by $n$ and $k$, subject to a condition local to a bounded window of consecutive
lines. Substituting the fixed index into the entry's own wording turns the definition into an
ordinary one-parameter array count, and for such a count the admissible windows are the
vertices of a finite digraph and an array is a walk. So the line is a walk count on $S$
vertices, and by the Cayley--Hamilton theorem it satisfies a monic integer recurrence of order
at most $S$.

\section{Recovering the recurrences}

\begin{lemma}\label{lem:bm}
A sequence known to satisfy some linear recurrence of order at most $S$ is determined, with its
minimal recurrence, by its first $2S$ terms; Berlekamp--Massey applied to them returns that
minimal recurrence exactly.
\end{lemma}

\subsection*{Column $k=3$}
Substituting the fixed index into the entry's wording gives an ordinary one-parameter array
count, a walk on $S=272$ vertices. Berlekamp--Massey on $2S=544$ exact terms
returns a minimal recurrence of order $10$ --- the order the entry states --- with
integer coefficients
\begin{center}\small\begin{tabular}{cccccccc}
$c_{1}$ & $10$ & $c_{4}$ & $-176$ & $c_{7}$ & $-76$ & $c_{10}$ & $-16$ \\
$c_{2}$ & $33$ & $c_{5}$ & $-300$ & $c_{8}$ & $79$ &  &  \\
$c_{3}$ & $6$ & $c_{6}$ & $-15$ & $c_{9}$ & $72$ &  &  \\
\end{tabular}\end{center}
so that $a(m)=\sum_{i=1}^{10}c_i\,a(m-i)$ for every $m$. Removing the first
$10+4$ terms and repeating gives order $10$ again, which is what the
identification below needs.

\subsection*{Column $k=4$}
Substituting the fixed index into the entry's wording gives an ordinary one-parameter array
count, a walk on $S=1056$ vertices. Berlekamp--Massey on $2S=2112$ exact terms
returns a minimal recurrence of order $30$ --- the order the entry states --- with
integer coefficients
\[
(c_1,\dots,c_{30})=(18,\ 162,\ 236,\ -2046,\ -9966,\ -6228,\ 65257,\ 188276,\ -46843,\ -649584,\ -764788,\ 412637,\ 1321769,\ 761732,\ 182925,\ -381900,\ -337601,\ -102296,\ 5318,\ -23355,\ -10567,\ -20176,\ -133,\ -430,\ -2201,\ 428,\ 116,\ 62,\ 8,\ 4).
\]
so that $a(m)=\sum_{i=1}^{30}c_i\,a(m-i)$ for every $m$. Removing the first
$30+4$ terms and repeating gives order $30$ again, which is what the
identification below needs.

\subsection*{Row $n=2$}
Substituting the fixed index into the entry's wording gives an ordinary one-parameter array
count, a walk on $S=72$ vertices. Berlekamp--Massey on $2S=144$ exact terms
returns a minimal recurrence of order $8$ --- the order the entry states --- with
integer coefficients
\begin{center}\small\begin{tabular}{cccccccc}
$c_{1}$ & $5$ & $c_{3}$ & $23$ & $c_{5}$ & $10$ & $c_{7}$ & $-3$ \\
$c_{2}$ & $11$ & $c_{4}$ & $45$ & $c_{6}$ & $1$ & $c_{8}$ & $-2$ \\
\end{tabular}\end{center}
so that $a(m)=\sum_{i=1}^{8}c_i\,a(m-i)$ for every $m$. Removing the first
$8+4$ terms and repeating gives order $8$ again, which is what the
identification below needs.

\subsection*{Row $n=3$}
Substituting the fixed index into the entry's wording gives an ordinary one-parameter array
count, a walk on $S=272$ vertices. Berlekamp--Massey on $2S=544$ exact terms
returns a minimal recurrence of order $20$ --- the order the entry states --- with
integer coefficients
\begin{center}\small\begin{tabular}{cccccccc}
$c_{1}$ & $9$ & $c_{6}$ & $-1668$ & $c_{11}$ & $4218$ & $c_{16}$ & $-152$ \\
$c_{2}$ & $51$ & $c_{7}$ & $-4457$ & $c_{12}$ & $4541$ & $c_{17}$ & $-224$ \\
$c_{3}$ & $222$ & $c_{8}$ & $-7039$ & $c_{13}$ & $9610$ & $c_{18}$ & $432$ \\
$c_{4}$ & $815$ & $c_{9}$ & $-6459$ & $c_{14}$ & $-285$ & $c_{19}$ & $512$ \\
$c_{5}$ & $770$ & $c_{10}$ & $2910$ & $c_{15}$ & $-4890$ & $c_{20}$ & $-256$ \\
\end{tabular}\end{center}
so that $a(m)=\sum_{i=1}^{20}c_i\,a(m-i)$ for every $m$. Removing the first
$20+4$ terms and repeating gives order $20$ again, which is what the
identification below needs.

\subsection*{Row $n=4$}
Substituting the fixed index into the entry's wording gives an ordinary one-parameter array
count, a walk on $S=1056$ vertices. Berlekamp--Massey on $2S=2112$ exact terms
returns a minimal recurrence of order $54$ --- the order the entry states --- with
integer coefficients
\[
(c_1,\dots,c_{54})=(17,\ 199,\ 2095,\ 17163,\ 49164,\ 5728,\ -982225,\ -9162832,\ -22727679,\ 22600243,\ 137134660,\ 472206655,\ 1184571797,\ -1171631903,\ -5466020520,\ -7174706705,\ -23109502971,\ 10347448009,\ 92285895608,\ 52003875308,\ 189224703657,\ 31476357249,\ -674091223527,\ -215277659729,\ -715597658065,\ -562065153852,\ 2168538834706,\ 323794079026,\ 1372177045551,\ 1690947992899,\ -3333605114781,\ -68242896950,\ -1400539218595,\ -1487594626433,\ 2268960028309,\ -181589577357,\ 611518034246,\ 462658501938,\ -623586266026,\ -461025094,\ -16432639110,\ -89621042484,\ 75680710988,\ 12159359476,\ -17846423736,\ 7876933224,\ -1393610496,\ -1790080416,\ 1186644672,\ -105946176,\ -44017920,\ 51867648,\ -11612160,\ -2239488).
\]
so that $a(m)=\sum_{i=1}^{54}c_i\,a(m-i)$ for every $m$. Removing the first
$54+4$ terms and repeating gives order $54$ again, which is what the
identification below needs.

\begin{theorem}
For each line above, the displayed recurrence holds for every index, and it is the recurrence
the entry records.
\end{theorem}

\begin{proof}
It holds by Lemma~\ref{lem:bm}. For the identification, the entry asserts that the line
satisfies SOME recurrence of the stated order, possibly only from some index on. Let $q$ be its
characteristic polynomial and $q_{\min}$ the minimal polynomial of the tail. Then
$q_{\min}\mid q$, and the second Berlekamp--Massey run --- on the line with its first few
terms removed --- shows $\deg q_{\min}$ equals the stated order, which is $\deg q$. Both are
monic, so $q=q_{\min}$.
\end{proof}

\section{Verification}

Three checks.

First, each line's model was compared against the entry's own published values for that line,
taken from the antidiagonal DATA read in either orientation and from the ``Table starts''
block, and agrees at every available term. This is the only point at which reading the entry's
English enters, and a misreading gives different counts.

Second, each recovered recurrence was evaluated directly on the model's values, with no
Berlekamp--Massey involved, and holds at every index tested.

Third, each order was computed twice by independent routes --- modulo a $61$-bit prime, where
every intermediate is a machine word, and again in exact rational arithmetic --- and once more
on the line with its first few terms dropped, which is what the identification requires. All
agree.

\begin{thebibliography}{9}
\bibitem{oeis} The OEIS Foundation, \emph{The On-Line Encyclopedia of Integer
Sequences}, \url{https://oeis.org}, 2026. Sequence A232076.
\bibitem{massey} J.~L.~Massey, Shift-register synthesis and BCH decoding, \emph{IEEE
Trans. Inform. Theory} \textbf{15} (1969), 122--127.
\bibitem{stanley} R.~P.~Stanley, \emph{Enumerative Combinatorics}, Volume 1, 2nd ed.,
Cambridge University Press, 2011. (Transfer-matrix method, Section 4.7.)
\bibitem{horn} R.~A.~Horn and C.~R.~Johnson, \emph{Matrix Analysis}, 2nd ed., Cambridge
University Press, 2013.
\end{thebibliography}

\end{document}
